\section{Conclusion}
\label{sec:conclusion}

Conditional inference forests are best understood as a method with a clean
calibrated core and a broader adaptive shell. Fixed-node Stage~A screening in
fixed-$B$ mode admits classical finite-sample permutation guarantees; the
larger tree and forest procedure does not inherit those guarantees
automatically.

The paper layout reflects that distinction. Theory is confined to the
fixed-node screening object, experiments are organized around a locked
rank-then-evaluate design, and broader empirical summaries are kept separate
from pre-specified confirmatory claims. This makes the manuscript easier to
update as the remaining results finalize, without changing the underlying
story.
